% =============================================================================
% Referee-revised patch: jet-Gram foundation
%   Replaces §2  sec:local-kernel-and-jet-normalization
%            §3  sec:jet-limit-local-blocks
%
% This version keeps the Gram-normalized convention, proves the theta-phase
% derivative estimates from differentiated Stirling asymptotics, states all
% regularity/separation hypotheses used in the Taylor limits, and separates
% the algebraic Gram criterion from the theta-specific positivity theorem.
% =============================================================================

% =============================================================================
\section{Local kernel and jet normalization}
\label{sec:local-kernel-and-jet-normalization}
\statuspaper{LIVE}\quad
\statusmig{migrated}\quad
\statusreferee{in-progress}\quad
\statussympy{verified}\quad
\statussim{partial}\quad
\statuslean{not-started}
\par\medskip

All local blocks in this paper are extracted from the Riemann--Siegel phase
on high-height real windows.  Throughout this section, a height parameter
\(T\) is assumed sufficiently large, and \(I_T\) denotes an interval centered
at \(T\) with
\[
        |I_T|\ll Q(T)^{-1},
        \qquad I_T\subset [T/2,2T].
\]
The second inclusion follows automatically at high height for every slowly
growing choice of \(Q(T)\) used below, but it is recorded because the
uniformity of the Stirling estimates is used on the whole window.

\subsection{The Riemann--Siegel phase chart}
\label{subsec:local-phase-chart}

\begin{definition}[Riemann--Siegel phase]
\label{def:riemann-siegel-phase}
For \(t>0\) let
\begin{equation}
\label{eq:theta-definition}
 \theta(t)
 := \Im \log \Gamma\!\left(\frac14+\frac{it}{2}\right)
    -\frac{t}{2}\log\pi,
\end{equation}
where the branch of \(\log\Gamma\) is the continuous branch obtained from
positive real argument.  Thus
\[
        \zeta\!\left(\frac12+it\right)
        =e^{-i\theta(t)}Z(t)
\]
with \(Z(t)\in\mathbb R\) for real \(t\).
\end{definition}

\begin{lemma}[Differentiated theta asymptotics]
\label{lem:theta-derivative-asymptotics}
Uniformly for \(t\in I_T\),
\begin{align}
\label{eq:theta-derivative-asymptotics}
q(t):=\theta'(t)
&=\frac12\log\frac{t}{2\pi}-\frac{1}{48t^2}+O(t^{-4}),\\
q'(t)=\theta''(t)
&=\frac{1}{2t}+O(t^{-3}),\\
q''(t)=\theta'''(t)
&=-\frac{1}{2t^2}+O(t^{-4}).
\end{align}
The constants in the \(O\)-terms are absolute.
\end{lemma}

\begin{proof}
Apply the differentiated Stirling expansion for \(\log\Gamma z\), uniformly
in the sector containing \(z=1/4+it/2\), and take imaginary parts in
\eqref{eq:theta-definition}; equivalently, differentiate the standard
Riemann--Siegel expansion
\[
\theta(t)=\frac{t}{2}\log\frac{t}{2\pi}-\frac{t}{2}-\frac{\pi}{8}
          +\frac{1}{48t}+O(t^{-3})
\]
with its differentiated remainder.  This gives the first line of
\eqref{eq:theta-derivative-asymptotics}; the remaining two lines follow by
one and two further differentiations of the differentiated Stirling expansion.
\end{proof}

\begin{definition}[Local phase chart]
\label{def:local-phase-chart}
A retained high-height local phase chart consists of the interval \(I_T\)
and the real smooth phase
\[
        \Phi_T(t):=\theta(t)\qquad (t\in I_T).
\]
We write \(\Phi\) and \(q=\Phi'\) when the center \(T\) is fixed.  The chart
conditions used downstream are:
\begin{enumerate}[label=\textup{(P\arabic*)},leftmargin=*]
\item \(\Phi_T=\theta\) on \(I_T\);
\item \(q(t)\ge 1\) on \(I_T\), after increasing the global lower-height
      cutoff if necessary.
\end{enumerate}
In particular, if \(Q(T)\ge1\), then (P2) gives the polynomial-in-\(Q\)
lower bound \(q(t)\ge Q(T)^{-c_q}\) with \(c_q=0\).
\end{definition}

\begin{lemma}[Phase-derivative lower bound]
\label{lem:phase-derivative-lower-bound}
Condition \textup{(P2)} of Definition~\ref{def:local-phase-chart} holds for
all sufficiently large \(T\).  More quantitatively, uniformly on \(I_T\),
\[
        q(t)\ge \frac12\log\frac{T}{4\pi}-C T^{-2}
\]
for an absolute constant \(C\).
\end{lemma}

\begin{proof}
By Lemma~\ref{lem:theta-derivative-asymptotics},
\(q(t)=\frac12\log(t/(2\pi))+O(t^{-2})\).  Since \(t\in I_T\subset[T/2,2T]\),
\(\log(t/(2\pi))\ge \log(T/(4\pi))\), and the stated bound follows.  The
right-hand side is eventually larger than \(1\).
\end{proof}

\begin{remark}[No zeta-side regularity is used here]
\label{rem:chart-regularity-zeta-side}
The present chart uses \(\zeta\) only through the identity
\(\zeta(1/2+it)=e^{-i\theta(t)}Z(t)\) and through the theta phase.  No lower
bound on \(|Z|\), no zero-free condition in \(I_T\), and no Hardy/source-energy
input is used in Sections~\ref{sec:local-kernel-and-jet-normalization}--
\ref{sec:jet-limit-local-blocks}.  Such conditions are to be imposed only in
the later sections where actual-zeta transverse suppression or source-energy
transfer is consumed.
\end{remark}

\subsection{The phase kernel}
\label{subsec:phase-kernel}

\begin{definition}[Phase kernel]
\label{def:phase-kernel}
For a real \(C^1\) phase \(\Phi\) on an interval \(I\), define
\begin{equation}
\label{eq:phase-kernel}
K_\Phi(x,y)=
\begin{cases}
\dfrac{\sin\bigl(\Phi(x)-\Phi(y)\bigr)}{\pi(x-y)}, & x\ne y,\\[1.2ex]
\dfrac{q(x)}{\pi}, & x=y,
\end{cases}
\qquad q=\Phi'.
\end{equation}
\end{definition}

\begin{lemma}[Basic kernel properties]
\label{lem:phase-kernel-properties}
The kernel \(K_\Phi\) is real-valued and symmetric.  If \(\Phi\in C^1\), its
extension across the diagonal is continuous and satisfies
\[
        K_\Phi(T,T)=\frac{q(T)}{\pi}.
\]
If \(\Phi\in C^4\), the diagonal first and mixed second derivatives exist and
are given by Lemma~\ref{lem:phase-kernel-diagonal-derivatives} below.
\end{lemma}

\begin{proof}
For \(x\ne y\), symmetry follows from
\(\sin(\Phi(y)-\Phi(x))/(y-x)=\sin(\Phi(x)-\Phi(y))/(x-y)\).  Real-valuedness
is immediate.  The diagonal value follows from
\(\Phi(y+h)-\Phi(y)=hq(y)+o(h)\).
\end{proof}

\begin{lemma}[Off-diagonal kernel derivatives]
\label{lem:phase-kernel-derivatives}
Let \(T_1\ne T_2\), set
\[
        s=T_1-T_2,
        \qquad \Delta=\Phi(T_1)-\Phi(T_2),
        \qquad q_j=q(T_j).
\]
Then, for \(\Phi\in C^1\) near \(T_1,T_2\),
\begin{align*}
K_\Phi(T_1,T_2) &= \frac{\sin\Delta}{\pi s},
&K_x(T_1,T_2) &= \frac{q_1s\cos\Delta-\sin\Delta}{\pi s^2},\\
K_y(T_1,T_2) &= \frac{\sin\Delta-q_2s\cos\Delta}{\pi s^2},
&K_{xy}(T_1,T_2)
&=\frac{(q_1+q_2)s\cos\Delta+(q_1q_2s^2-2)\sin\Delta}{\pi s^3}.
\end{align*}
\end{lemma}

\begin{proof}
Differentiate
\(K_\Phi(x,y)=\sin(\Phi(x)-\Phi(y))/\pi(x-y)\) by the quotient and chain
rules, then evaluate at \((T_1,T_2)\).
\end{proof}

\begin{lemma}[Diagonal kernel derivatives]
\label{lem:phase-kernel-diagonal-derivatives}
If \(\Phi\in C^4\) near \(T\), then
\begin{equation}
\label{eq:diagonal-kernel-derivatives}
K_\Phi(T,T)=\frac{q(T)}{\pi},\qquad
K_x(T,T)=K_y(T,T)=\frac{q'(T)}{2\pi},
\end{equation}
\[
        K_{xy}(T,T)=\frac{q''(T)+2q(T)^3}{6\pi}.
\]
\end{lemma}

\begin{proof}
Put \(x=T+u\), \(y=T+v\).  With derivatives evaluated at \(T\),
\[
\Phi(T+u)-\Phi(T+v)
=q(u-v)+\frac{q'}{2}(u^2-v^2)+\frac{q''}{6}(u^3-v^3)+O(|(u,v)|^4).
\]
Dividing by \(u-v\) and expanding \(\sin z=z-z^3/6+O(z^5)\) gives
\[
K_\Phi(T+u,T+v)
=\frac1\pi\left[
 q+\frac{q'}2(u+v)
 +\left(\frac{q''}{6}-\frac{q^3}{6}\right)(u^2+v^2)
 +\left(\frac{q''}{6}+\frac{q^3}{3}\right)uv
 +O(|(u,v)|^3)\right].
\]
The coefficients of \(u\), \(v\), and \(uv\) are the asserted derivatives.
\end{proof}

\subsection{Point-to-jet transform}
\label{subsec:point-to-jet-transform}

For symmetric pairs \(T\pm h\), define the Gram-normalized coalescing
even/odd transform
\begin{equation}
\label{eq:point-to-jet-transform}
P_h
=\frac1{\sqrt2}
\begin{pmatrix}
1&1\\[0.6ex]
-1/(2h)&1/(2h)
\end{pmatrix}
=\diag(\sqrt2,1/\sqrt2)
\begin{pmatrix}
1/2&1/2\\[0.6ex]
-1/(2h)&1/(2h)
\end{pmatrix}.
\end{equation}
Thus \(P_h\) is the literal central value/derivative jet transform followed
by the fixed diagonal congruence \(\diag(\sqrt2,1/\sqrt2)\).  We call this the
\emph{Gram-normalized} basis.

With this convention,
\begin{equation}
\label{eq:same-point-as-kernel-derivatives}
J(T)=
\begin{pmatrix}
2K_\Phi(T,T)&\partial_yK_\Phi(T,T)\\[0.6ex]
\partial_xK_\Phi(T,T)&\tfrac12\partial_{xy}K_\Phi(T,T)
\end{pmatrix},
\end{equation}
which is the same-point block computed explicitly in
Lemma~\ref{lem:same-point-jet-limit}.

\begin{remark}[Pair-separation convention]
\label{rem:pair-separation-convention}
The pair is \(T-h,T+h\), so its separation is \(2h\).  The lower row of
\(P_h\) therefore contains \(1/(2h)\).  Any formula imported from the
alternative \(T\pm h/2\) convention, or from the literal-jet matrix with a
lower row \((-1/h,1/h)\), must be re-conjugated before use.
\end{remark}

\begin{remark}[Convention audit]
\label{rem:normalization-audit}
With the literal central value/derivative transform at the same pair
\(T\pm h\), the same-point limit would be
\[
\frac1\pi
\begin{pmatrix}
q(T)&\tfrac12q'(T)\\[0.6ex]
\tfrac12q'(T)&\tfrac16\bigl(q''(T)+2q(T)^3\bigr)
\end{pmatrix}.
\]
The block used in this paper is obtained from that literal jet block by the
fixed diagonal congruence in \eqref{eq:point-to-jet-transform}; in particular
its \((1,1)\) entry is doubled and its \((2,2)\) entry is halved.  All
same-point blocks, cross-blocks, finite-scale blocks, whitening matrices, and
quotient covectors in the paper are to be read in this Gram-normalized
convention.
\end{remark}

% =============================================================================
\section{Jet-limit local blocks}
\label{sec:jet-limit-local-blocks}
\statuspaper{LIVE}\quad
\statusmig{migrated}\quad
\statusreferee{in-progress}\quad
\statussympy{verified}\quad
\statussim{partial}\quad
\statuslean{not-started}
\par\medskip

This section proves the two Taylor limits used downstream: the same-point
block \(J(T)\) and the cross-block \(N_{12}(T_1,T_2)\).  The final subsection
records positivity of \(J(T)\) for the actual theta phase.  For a general
phase, the algebraic determinant criterion remains valid, but positivity is
not automatic.

\subsection{Same-point jet limit}
\label{subsec:same-point-jet-limit}

\begin{lemma}[Same-point jet limit]
\label{lem:same-point-jet-limit}
Assume \(\Phi\in C^5\) near \(T\).  Let
\[
A_h(T)=
\begin{pmatrix}
K_\Phi(T-h,T-h)&K_\Phi(T-h,T+h)\\
K_\Phi(T+h,T-h)&K_\Phi(T+h,T+h)
\end{pmatrix},
\]
with row and column order \((T-h,T+h)\).  Then
\[
        P_hA_h(T)P_h^\top = J(T)+O(h^2)\qquad (h\to0),
\]
where
\begin{equation}
\label{eq:same-point-J}
J(T)=\frac1\pi
\begin{pmatrix}
2q(T)&\tfrac12q'(T)\\[1ex]
\tfrac12q'(T)&\tfrac1{12}\bigl(q''(T)+2q(T)^3\bigr)
\end{pmatrix}.
\end{equation}
\end{lemma}

\begin{proof}
All derivatives below are evaluated at \(T\).  Taylor expansion gives
\[
\Phi(T\pm h)=\Phi(T)\pm hq+\frac{h^2}{2}q'
             \pm\frac{h^3}{6}q''+\frac{h^4}{24}q'''
             +O(h^5),
\]
so
\[
\Phi(T-h)-\Phi(T+h)=-2hq-\frac{h^3}{3}q''+O(h^5).
\]
Using \(\sin z=z-z^3/6+O(z^5)\),
\begin{equation}
\label{eq:A12-same-point-expansion-revised}
K_\Phi(T-h,T+h)
=\frac{q}{\pi}+\frac{h^2}{6\pi}(q''-4q^3)+O(h^4).
\end{equation}
Also \(K_\Phi(T\pm h,T\pm h)=q(T\pm h)/\pi\).  For a symmetric matrix
\(M\), conjugation by \eqref{eq:point-to-jet-transform} gives
\begin{align*}
(P_hMP_h^\top)_{11}&=\tfrac12(M_{11}+2M_{12}+M_{22}),\\
(P_hMP_h^\top)_{12}&=\tfrac1{4h}(M_{22}-M_{11}),\\
(P_hMP_h^\top)_{22}&=\tfrac1{8h^2}(M_{11}-2M_{12}+M_{22}).
\end{align*}
Substitution of \eqref{eq:A12-same-point-expansion-revised} and the Taylor
series of \(q(T\pm h)\) yields
\begin{align*}
(P_hA_hP_h^\top)_{11}
&=\frac{2q}{\pi}+O(h^2),\\
(P_hA_hP_h^\top)_{12}
&=\frac{q'}{2\pi}+O(h^2),\\
(P_hA_hP_h^\top)_{22}
&=\frac{q''+2q^3}{12\pi}+O(h^2),
\end{align*}
which is \eqref{eq:same-point-J}.
\end{proof}

\subsection{Cross-block jet limit}
\label{subsec:cross-block-jet-limit}

\begin{lemma}[Cross-block jet limit]
\label{lem:cross-block-jet-limit}
Assume \(\Phi\in C^5\) in neighborhoods of \(T_1\) and \(T_2\).  Let
\(T_1\ne T_2\), set
\[
        s=T_1-T_2,
        \qquad \Delta=\Phi(T_1)-\Phi(T_2),
        \qquad q_j=q(T_j),
\]
and assume \(0<h<|s|/3\).  Define
\[
C_h(T_1,T_2)=
\begin{pmatrix}
K_\Phi(T_1-h,T_2-h)&K_\Phi(T_1-h,T_2+h)\\
K_\Phi(T_1+h,T_2-h)&K_\Phi(T_1+h,T_2+h)
\end{pmatrix},
\]
with row order \((T_1-h,T_1+h)\) and column order \((T_2-h,T_2+h)\).  Then
\[
        P_hC_h(T_1,T_2)P_h^\top
        =\frac1\pi N_{12}(T_1,T_2)+O(h^2)
        \qquad (h\to0),
\]
where
\begin{equation}
\label{eq:cross-N12}
N_{12}=
\begin{pmatrix}
\dfrac{2\sin\Delta}{s}
&
\dfrac{\sin\Delta-q_2s\cos\Delta}{s^2}
\\[2.4ex]
\dfrac{q_1s\cos\Delta-\sin\Delta}{s^2}
&
\dfrac{(q_1+q_2)s\cos\Delta+(q_1q_2s^2-2)\sin\Delta}{2s^3}
\end{pmatrix}.
\end{equation}
\end{lemma}

\begin{proof}
Since \(T_1\ne T_2\), the kernel is smooth near \((T_1,T_2)\).  Write
\(K_{ab}=\partial_x^a\partial_y^bK_\Phi(T_1,T_2)\).  Taylor expansion to total
order four gives, uniformly for \(u_1,u_2\in\{-h,+h\}\),
\[
K_\Phi(T_1+u_1,T_2+u_2)
=\sum_{a+b\le4}\frac{K_{ab}}{a!b!}u_1^a u_2^b+O(h^5).
\]
For a general \(2\times2\) matrix \(M\),
\begin{align*}
(P_hMP_h^\top)_{11}
&=\tfrac12(M_{11}+M_{12}+M_{21}+M_{22}),\\
(P_hMP_h^\top)_{12}
&=\tfrac1{4h}(-M_{11}+M_{12}-M_{21}+M_{22}),\\
(P_hMP_h^\top)_{21}
&=\tfrac1{4h}(-M_{11}-M_{12}+M_{21}+M_{22}),\\
(P_hMP_h^\top)_{22}
&=\tfrac1{8h^2}(M_{11}-M_{12}-M_{21}+M_{22}).
\end{align*}
Equivalently, these four expressions apply the sign weights
\(1\), \(\sigma_2\), \(\sigma_1\), and \(\sigma_1\sigma_2\), respectively,
to the four samples with \(u_i=\sigma_i h\).  Hence the parity of the Taylor
monomials selects
\[
(P_hC_hP_h^\top)_{11}=2K_{00}+O(h^2),\quad
(P_hC_hP_h^\top)_{12}=K_{01}+O(h^2),
\]
\[
(P_hC_hP_h^\top)_{21}=K_{10}+O(h^2),\quad
(P_hC_hP_h^\top)_{22}=\tfrac12K_{11}+O(h^2).
\]
Substituting the off-diagonal derivative formulas of
Lemma~\ref{lem:phase-kernel-derivatives} gives \eqref{eq:cross-N12}.
\end{proof}

\subsection{Same-point Gram positivity}
\label{subsec:same-point-gram-positivity}

\begin{definition}[Same-point Gram determinant]
\label{def:same-point-gram-determinant}
\begin{equation}
\label{eq:same-point-gram-determinant}
D_J(T):=4q(T)^4+2q(T)q''(T)-3q'(T)^2.
\end{equation}
\end{definition}

\begin{lemma}[Algebraic same-point Gram criterion]
\label{lem:algebraic-same-point-gram-criterion}
For the matrix \(J(T)\) in \eqref{eq:same-point-J},
\begin{equation}
\label{eq:J-trace-det}
\Tr J(T)=\frac{2q(T)^3+24q(T)+q''(T)}{12\pi},
\qquad
\det J(T)=\frac{D_J(T)}{12\pi^2}.
\end{equation}
Moreover,
\[
        J(T)\succ0
        \quad\Longleftrightarrow\quad
        q(T)>0\quad\text{and}\quad D_J(T)>0.
\]
\end{lemma}

\begin{proof}
The trace and determinant follow by direct computation from
\eqref{eq:same-point-J}.  Since \(J(T)\) is real symmetric, Sylvester's
criterion gives positive definiteness exactly when its leading principal minor
\(2q(T)/\pi\) and its determinant are positive.
\end{proof}

\begin{lemma}[Same-point Gram positivity for the theta phase]
\label{lem:same-point-gram-positivity}
For \(\Phi_T=\theta\), and for all sufficiently large retained heights,
\(J(T)\) is positive definite.  More precisely,
\begin{equation}
\label{eq:lambda-min-asymptotic}
        \lambda_{\min}(J(T))
        =\frac{2q(T)}{\pi}\,(1+o(1))
        =\frac1\pi\log\frac{T}{2\pi}\,(1+o(1)).
\end{equation}
In particular, after increasing the lower-height cutoff if necessary,
\(\lambda_{\min}(J(T))\ge1\); hence the polynomial lower bound
\(\lambda_{\min}(J(T))\ge Q(T)^{-C_J}\) holds with \(C_J=0\) whenever
\(Q(T)\ge1\).
\end{lemma}

\begin{proof}
By Lemma~\ref{lem:theta-derivative-asymptotics},
\[
q(T)=\frac12\log\frac{T}{2\pi}+O(T^{-2}),\qquad
q'(T)=O(T^{-1}),\qquad q''(T)=O(T^{-2}).
\]
Therefore
\[
D_J(T)=4q(T)^4+O\!\left(q(T)T^{-2}\right)+O(T^{-2})
      =4q(T)^4(1+o(1))>0
\]
for all sufficiently large \(T\), while \(q(T)>0\).  The algebraic criterion
of Lemma~\ref{lem:algebraic-same-point-gram-criterion} gives
\(J(T)\succ0\).

Let \(0<\lambda_1\le\lambda_2\) be the eigenvalues of \(J(T)\).  Since
\(\lambda_2\le\Tr J(T)\),
\[
\lambda_1=\frac{\det J(T)}{\lambda_2}
          \ge \frac{\det J(T)}{\Tr J(T)}
          =\frac{D_J(T)}{\pi(2q(T)^3+24q(T)+q''(T))}
          =\frac{2q(T)}{\pi}(1+o(1)).
\]
On the other hand, the Rayleigh quotient in the first coordinate gives
\(\lambda_1\le J_{11}(T)=2q(T)/\pi\).  This proves
\eqref{eq:lambda-min-asymptotic}.
\end{proof}

\begin{remark}[Scope of the positivity lemma]
\label{rem:scope-of-gram-positivity}
Lemma~\ref{lem:same-point-gram-positivity} is theta-specific.  For an arbitrary
real phase, the safe statement is only the algebraic criterion
\(q>0\) and \(D_J>0\).  In particular, this lemma does not prove any lower
bound for unrelated anisotropy functionals such as
\(q''-6q+2q^3\), nor does it supply Hardy/source-transfer information.
\end{remark}
